\documentclass[a4paper,11pt]{article} % Don't change this
\usepackage{fullpage} % Don't change this

\usepackage{titling} % Don't change this
\pretitle{\vspace{-6em}\begin{center}\larger[2]\bf} % Don't change this
\postauthor{\end{center}} % Don't change this
\preauthor{\vspace{-2em}\begin{center}\smaller[1]} % Don't change this
\postauthor{\end{center}\vspace{-5em}} % Don't change this

\usepackage[shortlabels]{enumitem} % Don't change this
\setlist{itemsep=0em,topsep=0.3em} % Don't change this

%%% Some useful packages.
\usepackage{xcolor}
\usepackage{float}
\usepackage{amsmath,amssymb,mathtools}
\usepackage{graphicx}
\usepackage[hidelinks]{hyperref}
\usepackage{relsize}
\usepackage{cleveref}
\crefname{table}{Table}{Tables}
\crefname{figure}{Figure}{Figures}

%%% Add additional packages or macros below
\usepackage{booktabs}
\usepackage{microtype}

% This file holds the text that brings the width arm (Tier 2) into the golden
% report. It is written against the golden report of 4 October 2026, the
% 21-page docs/report_golden.pdf, whose source is docs/report_golden.tex. The master source
% and its bibliography live elsewhere, so this file compiles on its own.
%
% Citations. \citepas{key}{text} prints "(text)" and \citetas{key}{text}
% prints "text". The key is the key of the golden report's ref.bib on the
% branch feat/grid. When the text is pasted into the master source, replace
% \citepas{key}{...} with \citep{key} and \citetas{key}{...} with \citet{key}.
% Every source cited here is already in the golden report's reference list.
%
% Cross references. \gref{label}{text} prints the number that the label has in
% the 21-page PDF. Replace it with \eqref{label} or \cref{label} when pasting.
% Where the pushed source has no label yet, the label is left empty.
\newcommand{\citepas}[2]{(#2)}
\newcommand{\citetas}[2]{#2}
\newcommand{\gref}[2]{#2}

\newcommand{\TODO}[1]{\textcolor{red}{\textbf{[TODO: #1]}}}
\newcommand{\where}[1]{\par\smallskip\noindent\textit{\small Where: #1}\par\smallskip}
\newcommand{\placeholderfigure}[1]{%
  \fbox{\parbox[c][4.5cm][c]{0.9\linewidth}{\centering\TODO{#1}}}}

\begin{document}

\title{The width arm: additions to the golden report}
\author{Group G2, Tier 2}
\date{}
\maketitle

\section*{How to use this file}

Each block below names the section of the golden report that it extends. It
also says whether its text is inserted as new text or replaces existing text.
Red markers stand for numbers and figures that wait on the width data. No
number from the Tier 2 runs of 26 September 2026 is used. Those runs share one
split across seeds, use the trace kernel and time grokking from memorisation,
so none of their numbers holds under the definitions of the golden report.

Section, equation, table and figure numbers refer to the 21-page golden report
of 4 October 2026.

% =====================================================================
\section*{Block 1. Section 1, Introduction}
\where{insert as a new paragraph after the paragraph that ends ``None of them
measures the shape of the kernel.''}

Width gives a test of these accounts that does not depend on the weight decay.
In \gref{eq:kernelflow}{(1)} the weight decay enters only the term
$-2(k-1)\lambda\Theta_t$. That term is proportional to the kernel, so it
changes the scale of the kernel and leaves its shape as it is. Only $T_t$ can
change the shape, and $T_0$ has expectation of order $n^{-1}$ in the width $n$
\citepas{lewkowyczTrainingDynamicsDeep2021}{Lewkowycz and Gur-Ari, 2021}. If
$T_t$ stays of that order along training and shape change sets the grokking
time, a wider network should grok later at a fixed weight decay. The grokking
time of \citetas{kimGrokkingWeightDecayClock2026}{Kim (2026)} scales as
$(1-\beta)/(\eta\lambda)$ and holds the width at most through a logarithmic
factor. We therefore add a width arm to the grid and ask which of the two the
event times follow.

% =====================================================================
\section*{Block 2. Section 3.1, Task, network and training}
\where{insert after the sentence that gives the readout factor
$c_N=N^{-1/2}$ of the NTK parameterisation and $c_N=N^{-1}$ of the mean-field
parameterisation.}

Every width in this report uses the NTK factor $c_N=N^{-1/2}$. In the
infinite-width limit the NTK stays constant and the network trains as a kernel
method \citepas{jacotNeuralTangentKernel2020}{Jacot, Gabriel, and Hongler,
2020}, so with this factor a wider network is lazier at a fixed $\alpha$. The
mean-field factor keeps the dynamics non-degenerate as the width grows
\citepas{chizatLazyTrainingDifferentiable2020}{Chizat, Oyallon, and Bach,
2020}, and under it the width would not act as a setting of laziness. A width
effect measured here belongs to the NTK parameterisation.

\where{replace the last sentence of the subsection, ``The grid sets the product
$\eta\lambda$ directly.''}

The grid sets the product $\eta\lambda$ and the width $N$ directly.

% =====================================================================
\section*{Block 3. Section 3.5, Event times}
\where{insert after the sentence that ends ``while (12) takes the pilot cell and
the cells of the second stage together, $\eta\lambda=0$ included, with
$\eta\lambda$ itself.''}

The width arm adds two fits of \gref{eq:aft}{(10)}. The first takes the cells
at the $\alpha$ of the pilot cell without weight decay, with the single
covariate $\ln N$. The second takes the cells with
$0<\eta\lambda\le\TODO{largest level of the crossed arm}$ at every width, with
the covariates $\ln N$, $\ln\eta\lambda$ and their product. Both logarithms are
centred at the middle of the grid before the product is formed, so the first
two coefficients are the slopes at that point. The width enters as $\ln N$.
Indicators of its levels would fail at the wide widths. When every seed of a
level is censored, the likelihood of an indicator for that level keeps rising
as its coefficient grows, so the indicator has no finite maximum likelihood
estimate. The model
\gref{eq:aalen}{(12)} takes $\ln N$ beside $\eta\lambda$. Its intervals
resample seeds with all their runs, as above, so the runs of one seed at every
width stay together. The relative risk model \gref{}{(13)} is also fitted with
$\ln N$ added to $Z(t)$. If the width acts on the event only through $S_t$,
$R_t$ and $A_t$, the coefficient of $\ln N$ in that fit should be close to zero
\citepas{kalbfleischStatisticalAnalysisFailure2002}{Kalbfleisch and Prentice,
2002}.

% =====================================================================
\section*{Block 4. Section 3.6, Choice of levels}
\where{insert as a new paragraph at the end of the subsection.}

The levels of the width arm were not chosen by this design. The widths form a
geometric grid with ratio 2,
$N\in\{50,100,200,400,800,1600\}$. The weight decays of the crossed arm,
$\eta\lambda\in\TODO{levels}$, lie below the lower end
$\eta\lambda=5\times10^{-4}$ of the scan. That is the range in which the
grokking time falls with the weight decay, and the fits of the second stage
say nothing about it. Every cell of the width arm has a budget of
$200\,000$ steps, twice the budget of the pilot cell and of the $\alpha$ arm.
A width cell is compared with those arms at the shorter budget, by treating
every event after $100\,000$ steps as censored there.

% =====================================================================
\section*{Block 5. New Section 4.2.3, Width}
\where{insert as a new subsubsection after Section 4.2.2, Weight decay.}

\subsubsection*{4.2.3\quad Width}

This subsection changes the width of the pilot cell and holds every other
setting fixed, with $\alpha=1$ and $\eta\lambda=0$. The widths are
$N\in\{50,100,200,400,800,1600\}$, each with \TODO{number of seeds} seeds and a
budget of $200\,000$ steps. \Cref{tab:width} gives the memorisation step and
the grokking step at level $0.95$ for each width, and \cref{fig:widthsurvival}
the Kaplan-Meier estimate of the probability of not yet having grokked.
\TODO{State how the memorisation step changes with $N$, and how the grokking
step changes with $N$. State at which widths no seed reaches each level within
the budget.} \Cref{tab:widthaft} gives the coefficient of $\ln N$ in
\gref{eq:aft}{(10)} for each event and level.

% Notes from the 5-seed dataset of 3 October 2026, kept here only as a guide to
% what to check. They use the sum kernel for A_t and must not be quoted. At
% alpha 1 and no weight decay the median memorisation step fell from about
% 10,000 steps at N = 50 to about 55 at N = 1600. Both N = 50 and N = 100
% grokked at level 0.95 on every seed. From N = 200 upwards no seed reached
% level 0.80 within 200,000 steps.

\begin{table}[H]
  \centering
  \caption{The sweep over width at the settings of the pilot cell. The
    memorisation level is training accuracy $0.99$ and the grokking level is
    test accuracy $0.95$. Medians, minima and maxima are over the seeds that
    reached the event.}
  \label{tab:width}
  \begin{tabular}{rrrrrr}
    \toprule
    & Memorised & \multicolumn{4}{c}{Grokked} \\
    \cmidrule(lr){2-2}\cmidrule(lr){3-6}
    $N$ & Median & Median & Minimum & Maximum & Reached \\
    \midrule
    50   & \TODO{} & \TODO{} & \TODO{} & \TODO{} & \TODO{} \\
    100  & \TODO{} & \TODO{} & \TODO{} & \TODO{} & \TODO{} \\
    200  & \TODO{} & \TODO{} & \TODO{} & \TODO{} & \TODO{} \\
    400  & \TODO{} & \TODO{} & \TODO{} & \TODO{} & \TODO{} \\
    800  & \TODO{} & \TODO{} & \TODO{} & \TODO{} & \TODO{} \\
    1600 & \TODO{} & \TODO{} & \TODO{} & \TODO{} & \TODO{} \\
    \bottomrule
  \end{tabular}
\end{table}

\begin{figure}[H]
  \centering
  \placeholderfigure{Kaplan-Meier curves by width at level 0.95, drawn as in
    Figure 5, with Greenwood's exponential interval dashed.}
  \caption{Kaplan-Meier estimates of the probability of not yet having grokked
    at level $0.95$, by width at the settings of the pilot cell, with
    Greenwood's exponential interval dashed
    \citepas{kalbfleischStatisticalAnalysisFailure2002}{Kalbfleisch and
    Prentice, 2002, Sec.~1.4.1}.}
  \label{fig:widthsurvival}
\end{figure}

\where{add these rows to Table 5, as the width arm.}

\begin{table}[H]
  \centering
  \caption{Rows of the width arm for Table 5: the coefficient of $\ln N$ in the
    log-normal accelerated failure time model \gref{eq:aft}{(10)}, with its
    $95\%$ interval, at the settings of the pilot cell.}
  \label{tab:widthaft}
  \begin{tabular}{llllrrr}
    \toprule
    Arm & Event & Level & Covariate & Coefficient & Lower & Upper \\
    \midrule
    width & memorised & 0.99 & $\ln N$ & \TODO{} & \TODO{} & \TODO{} \\
    width & grokked   & 0.80 & $\ln N$ & \TODO{} & \TODO{} & \TODO{} \\
    width & grokked   & 0.90 & $\ln N$ & \TODO{} & \TODO{} & \TODO{} \\
    width & grokked   & 0.95 & $\ln N$ & \TODO{} & \TODO{} & \TODO{} \\
    \bottomrule
  \end{tabular}
\end{table}

% =====================================================================
\section*{Block 6. New Section 4.2.4, Width and weight decay together}
\where{insert as a new subsubsection after Block 5.}

\subsubsection*{4.2.4\quad Width and weight decay together}

Section 4.2.2 finds that above the break more weight decay delays grokking,
and that its fits say nothing below $\eta\lambda=5\times10^{-4}$, where the
time falls. This subsection crosses that lower range with the widths of
Section 4.2.3, at $\alpha=1$, with $\eta\lambda\in\TODO{levels}$ and
\TODO{number of seeds} seeds per cell. \Cref{fig:widthdecay} gives the log
grokking time against $\ln\eta\lambda$ at each width, and \cref{tab:crossed}
the coefficients of the crossed fit of \gref{eq:aft}{(10)}. \TODO{State the
signs of the coefficients of $\ln N$ and $\ln\eta\lambda$ and whether each
interval excludes zero. State whether the product term is needed, and what its
sign says about how the effect of weight decay changes with width.}

\Cref{tab:widthphases} gives the phase of every run at the last step, at the
threshold $0.9$ of Section 4.2.2. \TODO{State the largest $\eta\lambda$ at
which each width still groks, and whether that edge moves as $N$ grows.}

% Notes from the 5-seed dataset of 3 October 2026, kept here only as a guide to
% what to check. At eta lambda = 3e-4 the runs at N = 800 memorised, but no
% seed reached level 0.95 within 200,000 steps. At N = 1600 no seed memorised.
% At eta lambda = 1e-3, 2 of the 5 seeds at N = 50 memorised and no seed at any
% other width did. At eta lambda = 3e-3 no seed memorised at any width.

The model \gref{eq:aalen}{(12)} with $\eta\lambda$ and $\ln N$ is fitted to
every cell of the width arm, $\eta\lambda=0$ included, and Figure 8 gains a
panel for $\ln N$. \TODO{State whether the cumulative regression function of
$\ln N$ rises, falls or stays flat along $t$.}

\begin{figure}[H]
  \centering
  \placeholderfigure{Log grokking time at level 0.95 against $\ln\eta\lambda$,
    one line per width. Filled marks are observed events and open marks are runs
    censored at the budget, as in Figure 4.}
  \caption{Log grokking time at level $0.95$ against $\ln\eta\lambda$ for each
    width of the crossed arm. Filled marks are observed events and open marks
    are runs censored at the budget.}
  \label{fig:widthdecay}
\end{figure}

\begin{table}[H]
  \centering
  \caption{Coefficients of the crossed fit of the log-normal accelerated
    failure time model \gref{eq:aft}{(10)} on the cells with $\eta\lambda>0$,
    with their $95\%$ intervals. The two logarithms are centred at the middle of
    the grid before their product is formed.}
  \label{tab:crossed}
  \begin{tabular}{lllrrr}
    \toprule
    Event & Level & Covariate & Coefficient & Lower & Upper \\
    \midrule
    memorised & 0.99 & $\ln N$                      & \TODO{} & \TODO{} & \TODO{} \\
    memorised & 0.99 & $\ln\eta\lambda$             & \TODO{} & \TODO{} & \TODO{} \\
    memorised & 0.99 & $\ln N\cdot\ln\eta\lambda$   & \TODO{} & \TODO{} & \TODO{} \\
    grokked   & 0.95 & $\ln N$                      & \TODO{} & \TODO{} & \TODO{} \\
    grokked   & 0.95 & $\ln\eta\lambda$             & \TODO{} & \TODO{} & \TODO{} \\
    grokked   & 0.95 & $\ln N\cdot\ln\eta\lambda$   & \TODO{} & \TODO{} & \TODO{} \\
    \bottomrule
  \end{tabular}

  \smallskip
  \TODO{add the rows for the grokking levels 0.80 and 0.90.}
\end{table}

\begin{table}[H]
  \centering
  \caption{Phase of the runs of the width arm at the last step, at the
    threshold $0.9$, following
    \citetas{pracherSpectralTheoryGrokking2026}{Pracher et al. (2026)}. Each
    entry gives the number of seeds that grok, memorise, forget or do not fit.}
  \label{tab:widthphases}
  \TODO{one row per width and one column per level of $\eta\lambda$,
    $\eta\lambda=0$ included.}
\end{table}

% =====================================================================
\section*{Block 7. Sections 4.3 and 4.4, the kernel across width}
\where{insert at the end of Section 4.3, The kernel at the grokking event.}

\Cref{tab:widthkernel} gives $S_t$, $R_t$ and $A_t$ at the grokking event at
level $0.95$ for each width without weight decay. $S_t$ and $R_t$ are taken on
the kernel of the summed logits and $A_t$ on the tangent kernel of all logits,
as in \gref{eq:statistics}{(6)}. \TODO{State how each statistic at the event
changes with $N$.}

\begin{table}[H]
  \centering
  \caption{The statistics \gref{eq:statistics}{(6)} at the grokking event at
    level $0.95$, by width at the settings of the pilot cell. Medians over the
    seeds that grokked, with the minimum and maximum in brackets.}
  \label{tab:widthkernel}
  \begin{tabular}{rlll}
    \toprule
    $N$ & $S_t$ & $R_t$ & $A_t$ \\
    \midrule
    50   & \TODO{} & \TODO{} & \TODO{} \\
    100  & \TODO{} & \TODO{} & \TODO{} \\
    200  & \TODO{} & \TODO{} & \TODO{} \\
    400  & \TODO{} & \TODO{} & \TODO{} \\
    800  & \TODO{} & \TODO{} & \TODO{} \\
    1600 & \TODO{} & \TODO{} & \TODO{} \\
    \bottomrule
  \end{tabular}
\end{table}

\where{add to Section 4.4, What times the event.}

Table 6 is fitted again on the runs of the width arm with $\ln N$ added to
$Z(t)$ (\cref{tab:widthcox}). \TODO{State whether the coefficient of $\ln N$ is
close to zero once $S_t$, $R_t$ and $A_t$ are in the model, and which of the
three carries the effect of width.}

\begin{table}[H]
  \centering
  \caption{Coefficients of the relative risk model with time-dependent
    covariates
    \citepas{kalbfleischStatisticalAnalysisFailure2002}{Kalbfleisch and
    Prentice, 2002, Eq.~6.14} on $S_t$, $R_t$, $A_t$ and $\ln N$, with their
    $95\%$ intervals and $p$-values, on the runs of the width arm.}
  \label{tab:widthcox}
  \begin{tabular}{lllrrrr}
    \toprule
    Event & Level & Covariate & Coefficient & Lower & Upper & $p$ \\
    \midrule
    grokked & 0.95 & $S_t$  & \TODO{} & \TODO{} & \TODO{} & \TODO{} \\
    grokked & 0.95 & $R_t$  & \TODO{} & \TODO{} & \TODO{} & \TODO{} \\
    grokked & 0.95 & $A_t$  & \TODO{} & \TODO{} & \TODO{} & \TODO{} \\
    grokked & 0.95 & $\ln N$ & \TODO{} & \TODO{} & \TODO{} & \TODO{} \\
    \bottomrule
  \end{tabular}

  \smallskip
  \TODO{add the rows for memorisation and for the grokking levels 0.80 and
    0.90.}
\end{table}

% =====================================================================
\section*{Block 8. Appendix B, Model checks}
\where{insert after Figure 7.}

\begin{figure}[H]
  \centering
  \placeholderfigure{Log-normal probability plot of the grokking times at level
    0.95 for each width, drawn as in Figure 7.}
  \caption{Log-normal probability plot of the grokking times at level $0.95$
    for each width of the width arm. Each observed time is placed at the normal
    quantile of its plotting position $(i-0.5)/n$, and runs censored at the
    budget are left out
    \citepas{nelsonAcceleratedTestingStatistical1990}{Nelson, 1990, Ch.~3,
    Sec.~4.2}. Parallel straight lines support the log-normal model
    \gref{eq:aft}{(10)}.}
  \label{fig:widthprobability}
\end{figure}

% =====================================================================
\section*{Decisions for the group}

These points depart from the golden report or are not settled by it. The
group should approve each one before the blocks above go into the master
source.

\begin{enumerate}
  \item The golden report has no width arm, and \texttt{CLAUDE.md} names a
    regression on width and weight decay as the head claim. The blocks above
    add width as a covariate of the report's own models. The group should
    confirm that the report takes this form.
  \item The width arm has a budget of $200\,000$ steps, where the pilot cell
    and the $\alpha$ arm have $100\,000$. Block 4 says how the arms are
    compared.
  \item The width arm has \TODO{number} seeds per cell, where the pilot cell
    and the $\alpha$ arm have 12.
  \item The width levels and the weight decays of the crossed arm were not
    chosen by the design of Section 3.6.
  \item The width enters the fits as $\ln N$ in place of indicators, for the
    reason given in Block 3.
  \item The crossed fit uses only the cells up to \TODO{largest level}. The
    cells above that level are kept in the Kaplan-Meier estimates, in the
    phases and in the model \gref{eq:aalen}{(12)}. The choice selects cells,
    and every run in a chosen cell is kept, so it is consistent with Section
    3.5. It must still be stated in the report.
  \item Section 4.2.3 reports level $0.95$. The golden report uses that level
    in every subsection except Section 4.2.2, which uses $0.90$.
\end{enumerate}

\end{document}
